\begin{description}
\item[Q1.] How much of PANSATZ's reported "few tens of iterations" convergence advantage vs a gate-based HEA is due to the pulse-duration parameterization itself, and how much is due to their steepest-ascent hill-climbing on a discrete grid, given that in my COBYLA reproduction the iteration count grew from $\sim$250 at short bond distances to the 500-cap at 2.5\,\AA?
\item[Q2.] What fraction of the residual energy error on real \texttt{ibm\_lagos} hardware (paper Fig.~5) is attributable to leakage-induced particle-number violation from the CR pulses, and would a Z-parity symmetry-verification post-selection recover a measurable fraction of that error, given that the two-qubit reduction is symmetry-based and PANSATZ explicitly allows leakage into higher transmon levels?
\item[Q3.] My ideal-simulator VQE error grew from $\sim 3\times 10^{-9}$\,Ha at 0.7\,\AA{} to $\sim 4\times 10^{-4}$\,Ha at 2.5\,\AA{}. Does PANSATZ's adaptive schedule (longer CR duration in the entangled regime) actually help the classical optimizer navigate this long-bond regime, or does it merely provide a more expressible circuit at fixed number of parameters? A Hessian-eigenvalue-spectrum comparison would disentangle these.
\item[Q4.] Is the paper's "PANSATZ 5 params vs GANSATZ 4 params" comparison for H$_2$ fair given that a 2-layer HEA ($\sim$12 params) already reaches nano-Hartree accuracy in my reproduction and would have comparable transpiled CR-pulse schedule duration?
\item[Q5.] Are the numerical Pauli-coefficient matrix elements of the paper's 2-qubit parity-reduced STO-3G H$_2$ Hamiltonian identical (down to sign conventions on the tapered qubit) to those produced by the current \texttt{qiskit-nature} 0.8 + \texttt{PySCF} 2.13 stack, so future PANSATZ ablations can be benchmarked against the paper's exact numbers rather than a numerically-close but subtly-different Hamiltonian?
\end{description}
